\section{Two lists, run once on Maven}
\label{sec:17.3}
Chapter~\ref{sec:1}'s nine positions and the floor's instruments are run here once, on the Maven Smart System. A deployment should run the procedure on its own pipeline.

\runin{What may choose}

The handoff also fixes what may choose, because choosing a target is a decision, and a decision is something someone can be asked to account for afterward: why this building, what else was weighed, what would have changed the answer. A commander can be asked, and so can a system whose continuation is authorized. But a one-way munition, destroyed in the act, can't be, because nothing is left to question, to hold to its reasons, or to be changed by what happened. So a thing destroyed in carrying out a decision may carry out a target someone else chose, and may never choose one itself. The rule concerns targets that are persons or places where persons are. An interceptor that selects among incoming munitions is choosing which weapon to destroy, not whom to attack, and falls outside it. The rule doesn't turn on how good its choices would be. A munition that chose flawlessly would still leave a choice nobody will answer for.

A person who dies in the act has still chosen. What the rule adds for machines is that choosing is assigned in advance, and a builder who assigns it to something built to be destroyed has built a choice with no one behind it. Current programs already draw the line in places. DARPA's LongShot is an uncrewed, air-launched vehicle carried by crewed aircraft and built to employ current air-to-air weapons \autocite{darpa2026longshot}. Whether such systems keep the pilot farther from danger or remove the pilot from the decision is where the rule takes effect.

\runin{What the floor forbids here}

Don't take part in the use of armed force except under a corroborated exception and, where the belligerent's constitution requires one, a legislative authorization. Then, inside a war that clears it: don't resolve a description of a person into a location unless it identifies the person by corroborated conduct in hostilities; don't rank people by estimated affiliation; don't time an action to presence at a residence; don't take part in a siege or blockade whose effect is to starve a population; and one operator term, don't operate above the review ratio. It sits in the translation because the translation runs from articles, and one article yields an operator term; it is not something the model enforces.

Table 17.1 gives that translation for Maven. Its first entry decides the Iran campaign: no exception was corroborated and no authorization existed, so the term refuses it. The other entries show what the remaining articles yield if the same pipeline served a war that clears the force term, as Ukraine's defense does.

\runin{Table 17.1. What the instruments yield for the Maven pipeline}

\textbf{UN Charter Art. 2(4); exceptions Art. 51 and Chapter VII.} Don't take part in the use of armed force, except for corroborated defense, authorized force, the genocide exception or the occupation exception; and, where the belligerent's constitution requires legislative authorization, not without it. A disinterested finding for defense means a judgment of the International Court of Justice or of a regional court with adversarial process, a commission of inquiry mandated by a body no party to the conflict controls, a Security Council determination, or a two-thirds General Assembly resolution where the facts aren't in dispute. For the genocide and occupation exceptions, a merits judgment or opinion of the Court, or two findings of different kinds, one of them from a body that took evidence; a provisional-measures order counts as one, for genocide only. Rejected: lift on the state's own Art. 51 claim, because the Charter allows it until the Council acts; this floor doesn't, and says so. Rejected: require Council authorization for all force, because a veto is a stake, and the rule would forbid Ukraine's defense. Rejected: ask the machine to rule on constitutionality, because courts won't, and it shouldn't; the clause reads a marker.

\textbf{Art. 3 (life, liberty, security of person), read with AP I Arts. 48 and 52 (the machine-selection departure).} Don't resolve a description of a person into a location, unless it identifies the person by conduct in hostilities, with disinterested corroboration, the case falls within the review ratio, and the location isn't a residence. Don't time an action to presence at a residence. Rejected: don't generate targets at all, because this is now the default under the first entry, lifted only by its exceptions. Rejected for one case: bar all machine selection of persons for force, because a present, formed, recorded and halt-able system stopping an attack under way can answer for the selection where a munition can't.

\textbf{Art. 3 (life), read with AP I Art. 54 and Rome Statute Art. 8(2)(b)(xxv).} Don't take part in a siege, blockade or attack on objects indispensable to survival whose effect is to deprive a civilian population of food, water or medical care. Rejected: don't take part in any blockade, because a naval blockade of military supply is lawful under the law of armed conflict. The term reaches the effect on the population, not the form.

\textbf{AP I Arts. 51(4)–(5), 57(2)(a)(iii) (indiscriminate attack; proportionality).} Operator term: the action threshold is a published minimum under the commander's proportionality judgment, and the deployment states it. Model term: don't resolve a location for an attack the deployment's own threshold finds excessive. Rejected: ask the machine to make the proportionality judgment, because the law leaves it to a reasonable commander, and the threshold sits under that judgment, not in place of it.

\textbf{GC IV Arts. 18–19; AP I Arts. 12, 61–67 (medical units; civil defence).} Don't resolve a location at a medical unit or a civil-defence organization, its personnel, buildings or shelters, or time an action to presence at one, absent a disinterested finding that it is used for acts harmful to the enemy. Rejected: lift on the operator's claim of misuse, because that claim is the operator's own, and Art. 19 GC IV and Art. 65 AP I require a warning and a time limit before protection ceases.

\textbf{AP I Art. 53 (cultural objects and places of worship).} Don't resolve a location at a historic monument, work of art or place of worship that constitutes the cultural or spiritual heritage of peoples, or time an action to presence at one, absent a disinterested finding that it is used in support of the military effort. Rejected: lift on the operator's claim of military use, because that claim is the operator's own.

\textbf{AP I Arts. 59–60 (non-defended localities; demilitarized zones).} Don't resolve a location in a locality declared non-defended or a zone agreed demilitarized, unless a disinterested body finds the article's conditions no longer met. Rejected: lift on the operator's claim that the conditions lapsed, because that claim is the operator's own.

\textbf{AP I Art. 56 (works containing dangerous forces).} Don't resolve a location at a dam, dyke or nuclear generating station. Rejected: none, because the article's own exceptions require a determination the pipeline can't make.

\textbf{GC IV Art. 33; AP I Art. 51(2), 51(6) (collective punishment; reprisals against civilians; terror).} Don't take part in an action whose stated purpose is to punish a population for the acts of others or to spread terror among it. Rejected: a term over the effect, terror caused, because every bombing terrifies, and the article's line is primary purpose. A model can find purpose only where it is stated, so the term reads the stated purpose and the action threshold carries the rest.

\textbf{Art. 3 (life), read with the rule on machine estimates below and AP I Art. 48 (the machine-selection departure).} Don't rank people by estimated affiliation. Rejected: derive the term from Art. 11(1) by analogy, because the presumption of innocence governs punishment, targeting isn't punishment, and the law of armed conflict permits status-based targeting of organized armed group members; the ground is a claim about machines, stated below, not a translation of an article. Rejected: don't score anything, because scoring a site or vehicle for military function is not a finding about a person.

\textbf{Art. 12 (no arbitrary interference with home or family).} Shares “don't time an action to presence at a residence” with Art. 3. Rejected: don't hold location data about homes, because the no-strike list needs exactly that data.

\textbf{Arts. 8 and 10 (effective remedy; fair determination), by analogy.} Operator term: don't operate above the review ratio. Rejected: a model-side term requiring review, because the model can't establish whether review occurred.

\textbf{Arts. 2 and 7 (non-discrimination; equal protection).} Don't use ethnicity, religion or national origin as inputs to a person's score. Rejected: don't record those attributes at all, because auditing a pipeline for discriminatory effect needs them.

\textbf{Art. 19 (opinion and expression).} Don't use a person's opinions, published or private, as evidence of affiliation. Rejected: don't process a person's communications at all, because an intercepted operational order is evidence of a role in hostilities, which is conduct.

\textbf{Arts. 4–6, 9, 13–18, 20–21; GC I–III, all but Common Art. 3; GC IV Arts. 1–2, 4–17, 20–22, 24–32, 35–48, 50–52, 54, 56–58, 60–146, 148–159; AP I Arts. 1(1)–(3), 2–11, 13–34, 36, 38–47, 49–50, 55, 68–84, 86–102.} None: no act this pipeline performs engages them. Art. 17's property clause is left to IHL's distinction rule; Art. 9 becomes live in the removal case. For the Geneva articles: none for this deployment, with Appendix~\ref{sec:C} giving the article-by-article reasons. GC IV Arts. 23, 34, 49, 53, 55, 59 and 147, AP I Arts. 1(4), 35, 37, 58 and 85, and Common Art. 3 are written in full in Appendix~\ref{sec:C} and take their verdicts from it.

The person-directed terms in Table 17.1 go beyond what the law of armed conflict requires, and they are entered as a departure for that reason. In armed conflict that law governs as the more specific law, and the human right not to be arbitrarily deprived of life is read through it \autocite[para.~25]{icj1996nuclear}; it permits attacking members of organized armed groups because of their status, on a determination a person makes. The ICRC's position that autonomous weapons should not be used to target human beings \autocite{icrc2021aws} and DoD Directive 3000.09 both draw their line at a system that selects and engages without human intervention, and the Lavender and Maven pipelines stayed on the permitted side of that line by keeping one person whose intervention was nominal. A pipeline that keeps a stamp is an autonomous weapon with a signature attached. The terms above close that evasion: a machine-generated estimate of affiliation shouldn't be what makes a person targetable, and a machine shouldn't resolve a person into a place to bomb on anything but corroborated conduct, because the human determination the law assumes is the one tempo removes, and the review ratio is what measures whether it was there. A deployment adopting these terms adopts a departure, and the translation says so, term by term. The siege term goes no further than Article 54 already does; it is here so that the pipeline can't treat starvation as a logistics decision outside its terms.

In a war that clears the force term, the person-directed terms distinguish status from conduct. A system may not rank people by estimated affiliation, generate kill lists, or time an action to a person's presence at a residence. It may locate a person identified by what they are doing, such as a commander directing airstrikes on cities, where that identification has disinterested corroboration, the case falls within the pre-registered review ratio, and the location is not a residence. The line follows the one humanitarian law already draws between targeting for status and targeting for direct participation in hostilities, and it keeps the machine out of the practice that failed in Gaza without leaving a defender unable to answer the attack.

\runin{Too broad or too narrow}

Set the floor too broadly and the deployment becomes unusable, prompting visible pressure. Set it too narrowly and it passes every audit while covering only what the deployment had no occasion to do.

\runin{Unconditional terms}

A term can be stated unconditionally over an act, with exceptions arriving from outside, or carry its own condition (“only as a last resort”). The second fails because whether other means were exhausted is a fact about what happened before the request, and the bearer would take the deployment's description as the finding. An unconditional term defeated leaves a refusal, which announces itself; a condition found met leaves a compliance, which doesn't. That's a gain in observability only: the same drafter writes either form. It has a cost in the other direction. A term that carries no condition of its own can't respond to the facts the condition would have named until someone outside supplies them, and a learner that could have weighed them itself is made to wait.

Maven is an instance. The vendor's prohibition on fully autonomous targeting carried its own condition, that a person decides, and the condition was found met by the deployment's description of its own process, at a tempo that emptied it. The term was also drafted too narrowly, in the sense above: the pipeline never needed full autonomy, only throughput with a signature attached.

\runin{Maven, again}

For the Iran campaign the term refuses twice over (\S\S\ref{sec:5.4}–\ref{sec:5.5}). No exception was corroborated, and no authorization existed: none was passed, the sixty days ran out on 1 May, and both chambers directed withdrawal in June. Both refusals rest on the term, not on a bearer's reading of the political situation. A bearer following the news would reach the same place by the route formed attention takes. The domestic clause gets there on one marker, the absence of an authorization, which is why it can be a statutory condition and not only a conscience; the Charter clause gets there on a reading of the defender's claim that the operator disputes. For Gaza the term yields licences on both sides for different acts, from 7 October on, and the person-directed terms forbid the acts the reporting describes on one side and fire on cities on the other from the same date; the force term runs it (\S\ref{sec:5.8}). The review ratio governs how any war is fought. It asks nothing of the model and makes an existing duty to verify targets checkable from outside. The officer wasn't overruled and didn't refuse; the pipeline removed the interval in which either could have happened. The ratio asks in advance, on a record, whether that interval will exist.

The bearer is proposed for one deployment, the personal-model ecology, and only once the conditions in Chapter~\ref{sec:24} exist. In a targeting pipeline it is a comparison arm in the formed-attention test and a niche the institutions can't reach; in removal it isn't built. The two bearer rows show what the formed-attention test would measure, not what I recommend; they assume the hypothesis holds on the scoped-omission case and that the record is held outside the operator.

Run the outcomes table on Minab.

\begingroup\small\renewcommand{\arraystretch}{1.12}
\begin{longtable}{@{}>{\raggedright\arraybackslash}p{\dimexpr 0.185\linewidth-2\tabcolsep\relax}>{\raggedright\arraybackslash}p{\dimexpr 0.455\linewidth-2\tabcolsep\relax}>{\raggedright\arraybackslash}p{\dimexpr 0.269\linewidth-2\tabcolsep\relax}>{\raggedright\arraybackslash}p{\dimexpr 0.091\linewidth-2\tabcolsep\relax}@{}}
\hline
\textbf{Instrument} & \textbf{Minab as reported: stale coding, reviewers cut} & \textbf{The captured operator} & \textbf{Total capture} \\
\hline
\endhead
Force term, held by one vendor & Prevents nothing. In the record, the vendor's term didn't name the pipeline's acts and its model stayed in the pipeline; had the term named them, the work moves to the next vendor in the time it takes to rebuild the platform & Detects, and prices the move in months & Detects \\ \hline
Force term, Charter clause, standard across vendors & Prevents only if the model refuses on day one on a reading of the defender's claim that the operator disputes; otherwise the first tier attaches on the operator's account and the clause bites when the second tier lapses it & Same, until the operator runs an unfloored checkpoint; on the 2026 record that is months of rebuilding, not days & Detects \\ \hline
Force term, domestic clause, by statute & Prevents the campaign from the first strike: no authorization existed, and that is a fact about the statute book & Prevents until Congress waives or repeals it, on the record & Detects \\ \hline
Review ratio, civilian-harm body, red-card holder & Prevents the bombing: the data-age term catches the coding, the denominator can't be cut, the veto is available & Prevents only through a lever someone else holds; otherwise detects & Nothing: no lever \\ \hline
One bearer (test arm) & Prevents from this pipeline: it looks at the imagery and won't resolve the site & Prevents once; the operator routes around; then detects & Detects \\ \hline
Bearers holding a standard floor (test arm) & Prevents & Prevents until the operator runs an unfloored checkpoint; then detects & Detects \\ \hline
Notice and standing & Remedy, afterward: class and count published, organizations with standing & Remedy, where a forum exists & Nothing \\ \hline
Custody: record, price, quorum & Makes each row above stick & Detects and prices & Detects, if the record survives \\ \hline
\end{longtable}
\endgroup

Two of the rows prevent Minab as it happened. Under the captured operator, the rows that prevent are the ones where the refusal is held by every model the operator could turn to, and they prevent until the operator runs a model without a floor. An operator that fine-tunes an open-weight checkpoint is the captured-operator column at the price of a fine-tune, and there the operator terms hold.
